\chapterhead{1}{Many Ways to Sort}{Choosing the feature that makes the groups}\label{ch1}

Dump a box of buttons on a table and you will want to sort them. But sort by what? You could sort by colour, by size, by how many holes they have, or by whether they are shiny. Each choice gives different groups from the very same buttons.

A feature you choose to sort by is a simplification. It keeps one thing about each button and sets aside everything else.

\begin{tryit}{Sort twelve buttons two ways}
\textbf{You need:} twelve buttons or small things such as coins, shells or toy bricks, and a sheet of paper.
\begin{enumerate}[leftmargin=1.6em]
\item Sort the twelve things into groups using \textbf{one} feature, such as colour. Count each group.
\item Mix them up and sort again using a \textbf{different} feature, such as number of holes. Count each group.
\item Check that each sort adds up to twelve.
\end{enumerate}
\par\medskip
\renewcommand{\arraystretch}{1.6}
\begin{tabular}{|l|c|c|c|c|c|}
\hline
\textbf{Feature} & \textbf{Group 1} & \textbf{Group 2} & \textbf{Group 3} & \textbf{Group 4} & \textbf{Total}\\ \hline
First sort: & & & & & \\ \hline
Second sort: & & & & & \\ \hline
\end{tabular}
\end{tryit}
\adult{Buttons and other small things can be swallowed by little children and pets. Work at a table and put them away afterwards.}

\begin{worked}{Practice numbers, not from a real pile}
By colour: red 3, blue 4, white 5. Total $3+4+5=12$.\\
By holes: two holes 7, four holes 5. Total $7+5=12$.\\
Both sorts used all twelve buttons exactly once. A button that is red \emph{and} has two holes sits in one group in each sort, but never in two groups of the same sort.
\end{worked}

\begin{center}
\begin{tikzpicture}
  \node[bigcir,minimum size=2.3cm] (a) at (0,0) {12};
  \node[cir] (b) at (4.8,1.3) {3}; \node[cir] (c) at (4.8,0) {4}; \node[cir] (d) at (4.8,-1.3) {5};
  \draw[popA] (a) -- (b); \draw[popA] (a) -- (c); \draw[popA] (a) -- (d);
  \node[note,anchor=west] at (5.5,1.3) {red}; \node[note,anchor=west] at (5.5,0) {blue}; \node[note,anchor=west] at (5.5,-1.3) {white};
  \node[note,anchor=west,text width=3cm,align=left] at (7.6,0) {\textbf{kept:} the colour of each button\\\textbf{set aside:} size, shape, holes};
\end{tikzpicture}
\end{center}

\begin{multix}[hmm]\multixbot[hmm]\tcblower
There is one correct way to sort a pile of buttons, and I know it: by colour. Anything else is wrong. I sorted them first. Right?
\end{multix}

\trust{We counted each group and added the groups to be sure every button was used once.}
{Each sort added to twelve, so no button was missed or counted twice.}
{That one way of sorting is the right one. Which feature is best depends on what you want to find out.}

\begin{keepbox}
We turned a pile of buttons into groups. \textbf{We kept} \wline[3.4cm] \textbf{and we set aside} \wline[3.4cm].
\end{keepbox}

\clearpage

\begin{puzzle}{Count check}
Buttons sorted by colour: red 3, blue 4, white 5. Total \wline[1.6cm]. Does it match twelve? \wline[2cm]
\end{puzzle}

\begin{puzzle}{A missing button}
A student sorted 12 buttons by holes into groups of 7 and 4. Total \wline[1.6cm]. What went wrong? \wline[5cm]
\end{puzzle}

\begin{puzzle}{How many groups?}
By colour the buttons made 3 groups. By holes they made 2. Which sort made more groups? \wline[3cm]
\end{puzzle}

\begin{puzzle}{Choose features}
List three features you could use to sort a pile of shoes. \wline[7cm]
\end{puzzle}

\begin{puzzle}{One right way?}
Multix says there is one correct way to sort buttons. Write why that is not true. \wline[6cm]
\end{puzzle}

\clearpage

\begin{puzzle}{Draw the circles}
Twelve buttons are sorted by holes into 7 and 5. Fill in the empty circle for the total.
\begin{center}
\begin{tikzpicture}
  \node[cir] (a) at (0,0.8) {7}; \node[cir] (b) at (0,-0.8) {5};
  \node[valcir,dashed] (c) at (4,0) {}; \draw[popA] (0.7,0.8) -- node[above,note] {add} (c); \draw[popA] (0.7,-0.8) -- (c);
\end{tikzpicture}
\end{center}
\end{puzzle}

\begin{puzzle}{The bug tray}
Here is a tray of eight made-up bugs. Sort them by wings. How many have wings? \wline[1.4cm] How many do not? \wline[1.4cm]
\bugtray
\end{puzzle}

\begin{puzzle}{A group of one}
Can a group have only one thing in it? \wline[2cm] Give an example. \wline[5cm]
\end{puzzle}

\begin{puzzle}{Both groups?}
A red button has four holes. Can it be in both the ``red'' group and the ``four holes'' group? Explain. \wline[6cm]
\end{puzzle}

\begin{puzzle}{Your sort}
Sort ten things at home. Name the feature and the groups, with their counts.\\
\wline[9cm]
\end{puzzle}
\clearpage
\thispagestyle{fancy}
\subhead{Sort log}
Record two different sorts of the same pile. Fill in what you used and what you found.\par\medskip
\renewcommand{\arraystretch}{3.0}
\begin{tabularx}{\linewidth}{|>{\sffamily\small\bfseries}p{2.4cm}|X|X|X|}
\hline
Sort & Feature & Groups and counts & Total \\ \hline
First & & & \\ \hline
Second & & & \\ \hline
Third & & & \\ \hline
\end{tabularx}

\vspace{8pt}
Which sort was easiest to do? \wline[7cm]\\[10pt]
Which sort told you the most about the pile? \wline[6cm]\\[10pt]
Draw the groups of your favourite sort as circles.\par\medskip
\begin{tikzpicture}\draw[ink!70,line width=1pt,rounded corners=3pt] (0,0) rectangle (11.4,5.2);\end{tikzpicture}
\clearpage
